\documentclass[11pt,a4paper]{article}
\usepackage[margin=2.2cm]{geometry}
\usepackage[french]{babel}
\usepackage{fontspec}
\setmainfont{Cambria}
\usepackage{amsmath,amssymb}
\usepackage[most]{tcolorbox}
\usepackage{enumitem}
\usepackage[hidelinks]{hyperref}

\definecolor{bleuenspd}{HTML}{1F4E79}
\definecolor{grisClair}{HTML}{F2F2F2}
\definecolor{vertOK}{HTML}{2E7D32}

\newtcolorbox{consignedur}{colback=grisClair, colframe=bleuenspd, boxrule=0.8pt, arc=2pt, top=3pt, bottom=3pt}
\newtcolorbox{donnees}{colback=white, colframe=bleuenspd!60, boxrule=0.6pt, arc=2pt, top=3pt, bottom=3pt}

\pagestyle{empty}

\begin{document}

\begin{center}
{\Large\bfseries Statistique Bayésienne --- EPD-SDIA-4112}\\[2mm]
{\LARGE\bfseries Examen de mi-parcours (Chapitres 1 et 2)}\\[2mm]
{\large Dr.\ Noulapeu Ngaffo Armielle}\\[2mm]
Master 1 Science des Données et IA \quad --- \quad Durée : \textbf{1 h 30} \quad --- \quad Documents non autorisés\\[1mm]
Barème : 20 points \quad (temps indicatif par exercice)
\end{center}

\vspace{2mm}
\begin{consignedur}
\small\textbf{Consignes.} Justifiez chaque calcul. Les résultats numériques seront donnés avec 3 à 4 chiffres significatifs. Les formules de conjugaison sont rappelées à la fin du sujet. \textbf{Barème :} Exercice 1 (4 pts, 15 min), Exercice 2 (5 pts, 20 min), Exercice 3 (6 pts, 25 min), Exercice 4 (5 pts, 25 min), dont 5 min de relecture.
\end{consignedur}

% ============================================================
\section*{Exercice 1 --- Théorème de Bayes : test de dépistage (4 pts, 15 min)}

Un test de dépistage a une sensibilité de \textbf{90~\%} et une spécificité de \textbf{96~\%}. La prévalence de la maladie dans la population testée est de \textbf{3~\%}.

\begin{enumerate}[label=\alph*)]
  \item Calculer la probabilité d'être malade sachant que le test est positif : $P(M \mid T^+)$. \hfill (2 pts)
  \item Si le test est négatif, quelle est la probabilité d'être quand même malade, $P(M \mid T^-)$ ? \hfill (1,5 pts)
  \item Interpréter ces deux résultats : que conclut-on sur la valeur du test positif et du test négatif ? \hfill (0,5 pt)
\end{enumerate}

% ============================================================
\section*{Exercice 2 --- Calcul simple de posterior binomial (4 pts, 15 min)}

On observe $y = 12$ succès sur $n = 25$ essais. A priori : $\theta \sim \mathrm{Beta}(3, 3)$.

\begin{enumerate}[label=\alph*)]
  \item Donner la loi a posteriori de $\theta$. \hfill (1 pt)
  \item Calculer l'espérance a posteriori $\mathbb{E}(\theta \mid y)$. \hfill (1 pt)
  \item En déduire l'estimateur de Bayes $\delta^\pi(y)$. \hfill (1 pt)
  \item Calculer la variance a posteriori. \hfill (1 pt)
\end{enumerate}

% ============================================================
\section*{Exercice 3 --- Comparaison de deux priors (6 pts, 25 min)}

Mêmes données que l'exercice 2 : $y = 12$, $n = 25$.

\begin{enumerate}[label=\alph*)]
  \item Prior $\mathrm{Beta}(1, 1)$ : donner la posterior et calculer son espérance. \hfill (1,5 pts)
  \item Prior $\mathrm{Beta}(5, 5)$ : donner la posterior et calculer son espérance. \hfill (1,5 pts)
  \item Quel prior donne le plus de poids aux données ? Justifier à partir de la taille effective du prior. \hfill (1,5 pt)
  \item Calculer l'estimateur de Bayes dans chaque cas et comparer les trois estimateurs (y compris celui de l'exercice 2). \hfill (1,5 pt)
\end{enumerate}

% ============================================================
\section*{Exercice 4 --- Beta-Binomial : taux de clics (5 pts, 20 min)}

Vous testez un nouveau bouton sur un site web. Sur 100 visiteurs, 15 cliquent.

\begin{enumerate}[label=\alph*)]
  \item Avec un prior $\mathrm{Beta}(1, 1)$, quelle est la distribution posterior du taux de clics ? \hfill (1 pt)
  \item Calculer la moyenne et l'intervalle de crédibilité à 95~\%. \hfill (1,5 pt)
  \item Avec un prior $\mathrm{Beta}(5, 20)$ (taux attendu 20~\%), quelle est la nouvelle posterior ? \hfill (1 pt)
  \item Comparer les deux approches : quel prior est le plus influencé par les données, et quel prior estime le mieux le taux observé ? \hfill (1,5 pt)
\end{enumerate}

% ============================================================
\vspace{3mm}
\begin{donnees}
\small\textbf{Rappels de cours.}
\begin{itemize}[leftmargin=4mm, itemsep=0pt]
  \item Théorème de Bayes : $p(\theta \mid D) \propto p(D \mid \theta)\, p(\theta)$.
  \item Beta--Binomial : si $\theta \sim \mathrm{Beta}(\alpha, \beta)$ et $y \mid \theta \sim \mathrm{Bin}(n, \theta)$, alors $\theta \mid y \sim \mathrm{Beta}(\alpha + y,\; \beta + n - y)$.
  \item $\mathbb{E}(\theta) = \dfrac{\alpha}{\alpha + \beta}$, \qquad $\mathrm{Var}(\theta) = \dfrac{\alpha\beta}{(\alpha+\beta)^2 (\alpha+\beta+1)}$.
\end{itemize}
\end{donnees}

\end{document}
